% =====================================================================
%  Chapter 15 — The Verdict, Part III
% =====================================================================

\chapter{What Would Change}
\label{ch:consequences}

This chapter closes the book by separating two things that the coverage of
this story has tended to run together: what a proof of the Riemann hypothesis
would change, and what is already true regardless.

\section{If the claim is confirmed}

Suppose the review of Chapter~14's Phase 3 succeeds. What follows?

\subsection{For number theory}

The consequences would be transformative and would take decades to absorb.

\begin{itemize}
\item \textbf{Error terms.} The prime number theorem would acquire the
$\operatorname{Li}(x) + O(x^{1/2}\log x)$ error term, and the Riemann explicit
formula would become an exact computational instrument rather than an
asymptotic one.
\item \textbf{Analytic number theory.} The whole apparatus built on
zero-free regions --- de la Vallée Poussin's method, Vinogradov's bounds,
the zero-density estimates --- would be re-based on a foundation that no
longer needs the $\exp(-c\sqrt{\log x})$ refinements, and a large part of the
subject's technical complexity would become unnecessary.
\item \textbf{L-functions.} For every Dirichlet $L$-function and, under GRH,
for the Dedekind zeta functions of number fields, the zero-free region would
be the critical line exactly. Explicit formulas for primes in arithmetic
progressions, in short intervals, and almost all of the analytic theory of
$L$-functions would sharpen correspondingly.
\item \textbf{Randomness.} The Katz--Sarnak\idx{Katz--Sarnak} programme, which proves
distributional results for primes \emph{under} RH, would become unconditional.
\item \textbf{Higher $L$-functions.} Under GRH, the same reflection argument
that this book traced in Chapter~7 --- non-vanishing to the right of the line
plus the functional equation gives zeros on the line --- would extend to
automorphic $L$-functions via the analogue of the functional equation, and
the Artin holomorphy\idx{Artin holomorphy conjecture} conjecture would become substantially more tractable.
\end{itemize}

\subsection{For mathematics as a subject}

The more interesting consequences are not about the Riemann hypothesis.

\begin{quote}
\textbf{A proof of RH by machine, reviewed by humans, would be the first
instance of a major mathematical theorem established in this manner.}
\end{quote}

The historical precedents are worth remembering. The four-colour theorem was
proved with a computer check in 1976, and it took decades for the community
to settle what that meant --- whether a computer-assisted proof is a proof,
and what standards apply to the checking. The eventual consensus was yes, with
conditions: the check must be independently reproducible, and the argument
must be written up in a form a reader can evaluate.

The same conditions apply here, and the machinery exists: the Lean FRO's
tools, \texttt{lean4checker}, the Palomar registry, the \texttt{Comparator}
protocol. What is missing is not a mechanism but a review.

\subsection{For AI and mathematical research}

The episode is also a substantial data point about AI in mathematics, and
the owner of this book's remit is as much about that as about the hypothesis.

The sequence of events is unusual and worth taking seriously in both
directions.

\begin{description}[style=nextline,leftmargin=0pt,font=\sffamily\bfseries\small,
  itemsep=5pt]
\item[\textbf{What it shows.}] A formal proof of the most famous open problem
in mathematics was, at least in outline, produced by a collaboration between a
specialist and a large language model, in a form that a compiler accepts. The
author's own account is that the human supplied the mathematical steering ---
finding the toy setting where two ideas combined --- and the model supplied
breadth and execution.

\item[\textbf{What it does not show.}] That the model found the proof. The
author's account is the opposite: the human recognised the key idea, and the
model was faster at execution than at the part the author could not follow.
And the claim has not been verified. A formal artifact that compiles is a
\emph{candidate} for a proof, not a proof, and the author's own decision to
publish only Lean has been correctly described by commentators as making the
claim unfalsifiable by discourse --- the code either checks or it does not ---
which is exactly why the statement review is the whole game.

\item[\textbf{The general lesson.}] The valuable output here was not a
finished theorem but a \emph{structured, checkable, 17-million-line candidate}
that a human could not have produced at all, and that a human with the right
expertise might still turn into a theorem. That is a different and more
modest claim than ``AI solved the Riemann hypothesis'', and it may be the
more important one. The August 2026 two-thirds theorem is the cleanest
evidence: it is a genuine, verified, independently reproduced mathematical
result of AI-assisted work, and it was not a proof of RH.
\end{description}

\section{If the claim is refuted}

Suppose instead that the review finds that the formalized statement is
weaker than advertised, or that an instance of \texttt{admitted} is assumed
rather than constructed.

That would be a genuinely valuable outcome, and it is worth being clear that
it would not be a vindication of scepticism. It would mean that a
175-thousand-file, axiom-free, sorry-free formalisation of a substantial
mathematical programme had a defect --- which is a normal, expected,
``Tuesday'' outcome in the formalisation literature, and which the mirror's
own status line anticipated.

It would also not mean that the Riemann hypothesis is still as far from
proved as it was in 2025. The August 2026 two-thirds theorem would remain a
genuine result. The formalisation would remain a substantial piece of work
that a community could build on. And the framework --- the
theta-lift-and-multiplicity programme, the Weil--Weil positivity engine, the
Rallis seesaw and Iwasawa apparatus, all of it formalised --- would remain
available.

\begin{quote}
This is worth saying because the history of the subject is full of cases where
a failed proof attempt left behind the tools for a later success. Bombieri's
proof of the Weil bound for function fields was preceded by years of
fruitless work on the error term. The theory of $L$-functions was built by
people who were trying and failing to prove the Riemann hypothesis.
\end{quote}

\section{What is true regardless}

The book's closing claim is this. Whatever happens to the claim, the following
remain true as of 21 November 2026, and they are the durable content of this
story.

\begin{enumerate}
\item \textbf{The Riemann hypothesis is not settled.} No proof is established.
It remains, by the most recent assessments, an open problem, and the
compilation of an artifact does not change that.

\item \textbf{More than two thirds of the zeros of $\zeta$ are simple and on
the critical line.} This is a theorem --- arXiv:2608.13637 --- proved with
the same machinery, formally verified in Lean, and independently reproduced by
a third party.

\item \textbf{The method works.} A theta-lift-based approach to Weil positivity
over the integers, producing rigorous results about the zeros of $\zeta$, is
now a demonstrated technique rather than a proposal. That is a genuine change
in the subject's toolkit.

\item \textbf{Formal verification at this scale is feasible.} 17.4 million
lines, an axiom-free axiom set, two independent compilations, and a
development that a specialist could plausibly review. The scale of the
artifact is not the obstacle to verification; the review is.

\item \textbf{The evaluation criteria are clear, and they are met by nobody
else yet.} Compile, check the axioms, replay the kernel, and then read the
statements. The first three have been done. The fourth has not.
\end{enumerate}

\section{A closing word about how to read announcements}

The most durable thing this book can offer a reader is not a verdict on one
claim. It is a procedure.

\begin{quote}
When a proof of something important arrives, and it arrives as code, ask five
questions in order.

\textbf{Does it compile, and did more than one person do it?}\idx{evaluation procedure} If not, nothing
else matters.

\textbf{What are its axioms?} If the list is not Lean's three standard ones,
stop.

\textbf{Is the statement the right one?} Not ``is it a statement about
$\zeta$'' --- is it \emph{Mathlib's} $\zeta$, and is it the full strength, or
a weakened variant that happens to be provable?

\textbf{What are the load-bearing inputs, and what are they, really?} Find the
five or six theorems that everything else depends on, and read their
\emph{statements} against the literature. Not their proofs --- their
statements.

\textbf{Does the author's account match the artifact?} Find the place where
the claim's own narrative and the code diverge, and resolve it.
\end{quote}

Applied to this case, the first two questions are answered, the third is
answered, the fourth is the open work, and the fifth is where the two
documents come apart. That is the state of the claim, and it is a perfectly
good thing for a reader to be able to say.

\begin{keytakeaways}

\item If confirmed, RH would transform the error terms in analytic number
theory and would be the first major theorem established in a machine-checked
form.
\item The AI dimension is best stated modestly: a structured, checkable
candidate was produced that a human could not have produced at all. The
August two-thirds theorem is the cleanest verified example.
\item If refuted, the outcome would still be valuable: the formalised
programme would remain as a toolkit, and the subject's history is full of
cases where a failed attempt supplied the tools for a later success.
\item What is true regardless: RH is not settled; more than two thirds of the
zeros are simple and on the line; the method works; verification at this
scale is feasible; and the evaluation procedure is known.
\end{keytakeaways}

\vspace{12pt}
\begin{quote}
\small\itshape
``It is natural to feel that this is an end, in my case a loss of a principal
guiding light, but what I saw during the course of progress so amazed me that
I believe it is one dawn across science of something far greater.''
\par
\vspace{6pt}
\hfill --- Levent Alpöge, \emph{Remark}, 2026
\end{quote}

\vspace{8pt}
\begin{quote}
\small\itshape
This book does not know whether the dawn has come. It has established what it
would take to find out, and it has established that the question is being
asked in a form that can be answered --- which is more than could be said for
almost every previous attempt at this problem.
\end{quote}
